% ============================================================
% RESOLVE Swiss — supplementary material
% Internal note, not for publication. Derivation of the power formula
% (equation (2) of the SAP). Build: pdflatex, bibtex, pdflatex x2.
% ============================================================
\documentclass[11pt,a4paper]{article}

\usepackage[utf8]{inputenc}
\usepackage[T1]{fontenc}
\usepackage{lmodern}
\usepackage[english]{babel}
\usepackage{microtype}
\usepackage[a4paper,margin=2.6cm]{geometry}
\usepackage{amsmath}
\usepackage[numbers,sort&compress]{natbib}
\usepackage[hidelinks]{hyperref}

\setlength{\parskip}{0.4em}
\setlength{\parindent}{0pt}

\title{\large\bfseries Supplementary material to the RESOLVE Swiss
statistical analysis plan:\\ derivation of the power formula}
\author{}
\date{}

\begin{document}
\maketitle

The statistical analysis plan evaluates power with
\begin{equation}
\label{eq:power}
\mathrm{power} = 1 - T_{\nu,\lambda}\!\left(t_{1-\alpha/2,\,\nu}\right)
               + T_{\nu,\lambda}\!\left(-t_{1-\alpha/2,\,\nu}\right),
\qquad
\lambda = \frac{\delta}{\sqrt{\sigma^2 (1-r^2)\,\mathrm{DE}
\left(\frac{1}{\bar m k_1}+\frac{1}{\bar m k_2}\right)}},
\end{equation}
where $\delta$ is the target difference, $\sigma$ the between-participant
standard deviation of the outcome, $r$ the baseline-to-follow-up
correlation, $\mathrm{DE}$ the design effect, $\bar m$ the mean number of
analysed participants per practice, $k_1$ and $k_2$ the numbers of practices
per arm, $\nu = k_1 + k_2 - 2$, and $T_{\nu,\lambda}$ the distribution
function of the noncentral $t$ distribution with $\nu$ degrees of freedom
and noncentrality $\lambda$. This note derives the formula in five
elementary steps. Each step is standard, and the formula is their
composition.

\subsection*{Step 1: the estimator and its variance}

The planning test compares the two arm means, $\bar y_1 - \bar y_0$. In one
arm with $k$ practices of average size $\bar m$, clustering inflates the
variance of the arm mean by the design effect,
\begin{equation}
\operatorname{Var}(\bar y) = \frac{\sigma^2\,\mathrm{DE}}{\bar m k},
\qquad
\mathrm{DE} = 1 + \left[ \left( \mathrm{cv}^2\,\tfrac{k-1}{k} + 1 \right)
\bar m - 1 \right] \rho ,
\end{equation}
with $\rho$ the intra-cluster correlation and $\mathrm{cv}$ the coefficient
of variation of practice size \citep[eq.~1]{eldridge2006}. The two arms are
independent, so their variances add:
\begin{equation}
\label{eq:vardiff}
\operatorname{Var}(\bar y_1 - \bar y_0)
= \sigma^2\,\mathrm{DE} \left( \frac{1}{\bar m k_1} + \frac{1}{\bar m k_2}
\right).
\end{equation}
This is the variance behind the sample size formula of
\citet[eq.~7.12]{hayes2017}, with the per-arm terms added separately for
unequal numbers of clusters as in their section~7.6.2.

\subsection*{Step 2: the baseline}

The analysis uses the baseline measurement of the outcome, which multiplies
the relevant residual variance by the factor $1-r^2$ \citep{borm2007}. The
true standard error of the comparison is therefore
\begin{equation}
\label{eq:se}
\mathrm{SE} = \sqrt{\sigma^2 (1-r^2)\,\mathrm{DE}
\left( \frac{1}{\bar m k_1} + \frac{1}{\bar m k_2} \right)},
\end{equation}
the expression under the root in the noncentrality parameter of
eq.~\eqref{eq:power}.

\subsection*{Step 3: if the standard error were known}

With known $\mathrm{SE}$ the statistic $T = (\bar y_1 - \bar
y_0)/\mathrm{SE}$ would be standard normal under the null hypothesis and
normal with mean $\lambda = \delta/\mathrm{SE}$ and unit variance under a
true difference of $\delta$. Power would be $\Phi(\lambda - z_{1-\alpha/2})$.
This is the normal approximation, and it is the basis of the requirement
table in the analysis plan.

\subsection*{Step 4: the standard error is estimated from few practices}

In a cluster randomised trial the scarce ingredient of
$\widehat{\mathrm{SE}}$ is the between-practice variance. It is estimated
from the $k_1 + k_2$ practice-level summaries, and their pooled variance
carries $\nu = k_1 + k_2 - 2$ degrees of freedom
\citep[section~10.3.1]{hayes2017}, so that $s^2 \sim \sigma^2_{\mathrm{SE}}
\chi^2_\nu / \nu$ independently of the numerator. A ratio of the form
\begin{equation}
T = \frac{Z + \lambda}{\sqrt{\chi^2_\nu / \nu}},
\qquad Z \sim N(0,1) \text{ independent of } \chi^2_\nu,
\end{equation}
is by definition noncentral $t$ with $\nu$ degrees of freedom and
noncentrality $\lambda$. The test statistic has exactly this structure: the
numerator is the estimate, which in units of $\mathrm{SE}$ is a standard
normal shifted by $\lambda$, and the denominator is the estimated standard
error divided by the true one. Under the null hypothesis $\lambda = 0$, the
distribution is central $t_\nu$, and the test rejects when $|T| >
t_{1-\alpha/2,\,\nu}$.

\subsection*{Step 5: power is the rejection probability under $\lambda$}

\begin{equation}
\mathrm{power}
= P\!\left(|T| > t_{1-\alpha/2,\,\nu}\right)
= 1 - T_{\nu,\lambda}\!\left(t_{1-\alpha/2,\,\nu}\right)
    + T_{\nu,\lambda}\!\left(-t_{1-\alpha/2,\,\nu}\right),
\end{equation}
which is eq.~\eqref{eq:power}. This last step is the two-sample power
computation of \citet[section~3.2.1]{chow2018}, applied at the practice
level: their degrees of freedom $n_1 + n_2 - 2$ for individuals become
$k_1 + k_2 - 2$ for practices, and their standard error becomes
eq.~\eqref{eq:se}.

\subsection*{Assumptions}

Three approximations are inherited and are declared in the analysis plan.
The practice-level summaries are treated as normally distributed with a
common variance, which is the assumption of the underlying $t$ comparison
\citep[section~10.3.1]{hayes2017}. The design effect with unequal practice
sizes is the approximation of \citet[eq.~1]{eldridge2006}, exact for the
size-weighted comparison of practice means and conservative for
individual-level analyses. The baseline factor $1-r^2$ is applied at the
individual level, where a formula carrying the individual-level and the
cluster-level baseline correlations separately exists
\citep{teerenstra2012}. Evaluating power on the $t$ scale rather than the
normal scale builds the small-sample correction into the planning, as
current guidance recommends \citep{leyrat2018, vanbreukelen2018}.

The script \texttt{R/sample\_size.R} in the repository implements
eq.~\eqref{eq:power} line by line, and the simulations in
\texttt{R/simulations/} confirm it empirically.

\bibliographystyle{unsrtnat}
\bibliography{../references}

\end{document}
